% Scattering Angle in the CM System
%
\subsubsection{Change of variables}
%%%%%%%%%%%%%%%%%%%%%%%%%%%%%%%%%%%
%
The integral in \cref{eq:sn_reduced_theta} is of the form
%
\begin{equation}
    \label{eq:sn_legacy_integration_integral_R}
    \int_{R_0}^\infty F(R) \, \mathrm{d}R .
\end{equation}
%
where \(R_0\) is a pole of $F(R)$, which is obtained by solving \cref{eq:sn_apsis}, see Section~\ref{s:apsis}.
Having a numerically determined pole of the integrand as an integration boundary is inconvenient, so the first step of calculating the scattering integrals is to apply a change of variable.
One possibility is \cite{Eck91}
%
\begin{equation}
    \label{eq:sn_legacy_transform1}
    R = \frac{ R_0 }{ (1-u^2) } .
\end{equation}
%
yielding
%
\begin{equation}
    \label{eq:sn_legacy_integration_integral_u1}
    \int_{R_0}^\infty F(R) \, \mathrm{d}R =
        2 R_0 \int_0^1 F \left( \frac{ R_0 }{ 1-u^2 } \right)
        \frac{ u \, \mathrm{d}u }{ (1-u^2)^2 } .
\end{equation}
%
With \cref{eq:sn_reduced_theta} we obtain
%
\begin{equation}
    \label{eq:sn_legacy_cmangl_legendre_theta_u}
    \theta = \pi - \frac{ 4 P }{ R_0 } \int_0^1
        \frac{ u \, \mathrm{d}u }{ g \left( \frac{ R_0 }{ 1-u^2 } \right) } .
\end{equation}
%
where $g(R)$ is given by \cref{eq:sn_reduced_g}.

\subsubsection{The integrand at the limits of integration}
%%%%%%%%%%%%%%%%%%%%%%%%%%%%%%%%%%%%%%%%%%%%%%%%%%%%%%%%%%
%
First, we show that the integrand is well-behaved.
For that, we write \(g(R_0/(1-u^2))\) explicitly
%
\begin{equation}
    \label{eq:sn_legacy_cmangl_legendre_g_u}
    g \left( \frac{ R_0 }{ 1-u^2 } \right) =
    \sqrt{ 1 - \frac{ 1 }{ \varepsilon R_0 }
        \Phi\left( \frac{ R_0 }{ 1-u^2 } \right) (1-u^2)
        - \frac{ P^2 }{ R_0^2 } (1-u^2)^2 } .
\end{equation}
%
By subtracting \cref{eq:sn_apsis} from the expression under the square root, \cref{eq:sn_legacy_cmangl_legendre_g_u} may be rewritten
%
\begin{equation}
    \label{eq:sn_legacy_cmangl_legendre_g_chi}
    g\left( \frac{ R_0 }{ 1-u^2 } \right) =
        \frac{ u }{ \sqrt{ \varepsilon R_0 } } \sqrt{ \chi(u) + \rho (2-u^2) }
\end{equation}
%
with
%
\begin{equation}
    \label{eq:sn_legacy_cmangl_legendre_chi}
    \chi(u) = \frac{ \Phi(R_0) - \Phi\left( \frac{ R_0 }{ 1-u^2 } \right) \, (1-u^2) }{ u^2 }
\end{equation}
%
and
%
\begin{equation}
    \label{eq:sn_legacy_cmangl_legendre_rho}
    \rho = \frac{ \varepsilon P^2 }{ R_0 } .
\end{equation}

\(\chi(u)\) is well-behaved with
%
\begin{equation}
    \label{eq:sn_legacy_cmangl_legendre_chi_lim}
    \chi(1) = \Phi(R_0), \qquad
        \lim_{u \rightarrow 0} \chi(u) = \Phi(R_0) - R_0 \Phi'(R_0) .
\end{equation}
%
The latter is obtained using l’Hospital’s rule.
\(\chi(u)\) is plotted in \cref{fig:sn-cmangl-legendre-chi-plot} for the universal ZBL potential and some values of \(R_0\).
%
\begin{figure}[htbp]
    \centering
    \includegraphics[scale=0.6]{chi_ZBL.pdf}
    \caption{Function \(\chi(u)\) according to \cref{eq:sn_legacy_cmangl_legendre_chi} for the ZBL universal potential.}
    \label{fig:sn-cmangl-legendre-chi-plot}
\end{figure}

\Cref{eq:sn_legacy_cmangl_legendre_theta_u} may now be rewritten
%
\begin{equation}
    \label{eq:sn_legacy_cmangl_legendre_theta_chi}
    \theta = \pi - 4 \int_0^1
    \sqrt{ \frac{ \rho }{ \rho \, (2-u^2) + \chi(u) } }
    \, \mathrm{d}u .
\end{equation}
%
Since \(\chi(u)\) is well-behaved, so is the term under the square root.
The integrand of \cref{eq:sn_legacy_cmangl_legendre_theta_chi} therefore has no poles, and the integral is amenable to standard methods of numerical quadrature.

\subsubsection{Integration using Gauss-Legendre quadrature}
%%%%%%%%%%%%%%%%%%%%%%%%%%%%%%%%%%%%%%%%%%%%%%%%%%%%%%%%%%%
%
Using Gauss-Legendre quadrature, we can approximate \cref{eq:sn_legacy_cmangl_legendre_theta_chi} by
%
\begin{equation}
    \label{eq:sn_legacy_cmangl_legendre_theta_chi_approx_plain}
    \theta \approx
        \pi - 4 \sum_{i=1}^n \tilde{w}_i
        \sqrt{ \frac{ \rho }{ \rho \, (2-\tilde{u}_i^2) + \chi(\tilde{u}_i) } }
\end{equation}
%
with \(\tilde{u}_i\) and \(\tilde{w}_i\) defined in \cref{eq:sn_legacy_legendre_transform} of Appendix~\ref{s:gauss-legendre}.

\Cref{eq:sn_legacy_cmangl_legendre_theta_chi_approx_plain}, however, has the problem that it does not tend to zero exactly as \(P \rightarrow \infty\).
Note that for \(\rho \rightarrow \infty\) the integral in \cref{eq:sn_legacy_cmangl_legendre_theta_chi} becomes
%
\begin{equation}
    \label{eq:sn_legacy_cmangl_legendre_pi}
    \int_0^1 \frac{ 1 }{ \sqrt{ 2-u^2 } } \, \mathrm{d}u = \frac{ \pi }{ 4 } ,
\end{equation}
%
and \(\theta \rightarrow 0\) as it should.
However, since the integrand is not a polynomial, the Legendre approximation of the integral will not evaluate exactly to \(\pi/4\), and \cref{eq:sn_legacy_cmangl_legendre_theta_chi_approx_plain} will not tend to zero as \(\rho \rightarrow \infty\).
This may be fixed by scaling the integral to the correct value of \(\pi/4\):
%
\begin{equation}
    \label{eq:sn_legacy_cmangl_legendre_theta_chi_approx}
    \theta \approx
        \pi \left( 1
        - \frac{ 4 }{ \pi_\mathrm{num} }
        \sum_{i=1}^n \tilde{w}_i \sqrt{
            \frac{ \rho
            }{ \rho \, (2-\tilde{u}_i^2) + \chi(\tilde{u}_i) } }
        \right)
\end{equation}
%
with
%
\begin{equation}
    \label{eq:sn_legacy_cmangl_legendre_pinum}
    \pi_\mathrm{num} =
        4 \sum_{i=1}^n \tilde{w}_i \frac{ 1 }{ \sqrt{ (2-\tilde{u}_i^2) } }
\end{equation}
%
Since the sum in \cref{eq:sn_legacy_cmangl_legendre_theta_chi_approx} evaluates to \(\pi_\mathrm{num}/4\) for \(\rho \rightarrow \infty\), \(\theta \rightarrow 0\).
At the same time, for \(\rho \rightarrow 0\) (i.e., \(P \rightarrow 0\)) the sum goes to zero and \(\theta \rightarrow \pi\) as it should.

One can also apply this correction to the Legendre approximation of \cref{eq:sn_legacy_cmangl_legendre_theta_u}:
%
\begin{equation}
    \label{eq:sn_legacy_cmangl_legendre_theta_g_approx}
    \theta \approx
        \pi \left( 1 - \frac{ 4 }{ \pi_\mathrm{num} } \frac{ P }{ R_0 }
        \sum_{i=1}^n
        \frac{ \tilde{w}_i \tilde{u}_i
        }{ g \left( \frac{ R_0 }{ 1-\tilde{u}_i^2 } \right) } \right).
\end{equation}
%
\Cref{eq:sn_legacy_cmangl_legendre_theta_g_approx} is not numerically equivalent to \cref{eq:sn_legacy_cmangl_legendre_theta_chi_approx}.
It does not require the evaluation of \(\Phi(R_0)\), which might represent a non-negligible saving of computation time if few abscissas are used.
For few abscissas, the results of \cref{eq:sn_legacy_cmangl_legendre_theta_g_approx} and \cref{eq:sn_legacy_cmangl_legendre_theta_chi_approx} are very similar.
With increasing number of abscissas, the result is increasingly dependent on the accuracy of the apsis \(R_0\).
The term under the square root may even become negative, if the accuracy of \(R_0\) is low and the number of abscissas is large.

The integrand of \cref{eq:sn_legacy_cmangl_legendre_theta_chi} is shown in \cref{fig:integrand-legendre-theta} for the limiting cases \(\rho \gg \Phi, \chi\) and \(\rho \ll \Phi, \chi\).
It is smooth, although not completely polynomial (note the long flat parts of most curves).%
%
\footnote{Alternatively, one may insert \cref{eq:sn_legacy_cmangl_legendre_pinum} in \cref{eq:sn_legacy_cmangl_legendre_theta_chi_approx_plain} instead of \(\pi\), and one gets after some manipulation and correction with \(\pi / \pi_\mathrm{num}\)
%
\begin{equation*}
    \theta \approx 4 \frac{ \pi }{ \pi_\mathrm{num} } \sum_{i=1}^n
        \tilde{w}_i
        \frac{ \chi(\tilde{u}_i)
        }{ \sqrt{ 2-\tilde{u}_i^2 } \,
            ( S_1(\tilde{u}_i) + S_2(\tilde{u}_i) ) \,
            S_2(\tilde{u}_i) } .
\end{equation*}
%
with
%
\begin{equation*}
    S_1(u) = \sqrt{ \rho \, (2-u^2) }, \qquad
    S_2(u) = \sqrt{ \rho \, (2-u^2) + \chi(u) } .
\end{equation*}
%
This also guarantees \(\theta \rightarrow 0\) as \(\rho \rightarrow \infty\), and due to the correction also \(\theta \rightarrow \pi\) as \(\rho \rightarrow 0\).
While this has, in principle, better roundoff behavior, the improvement is insignificant for the small numbers of abscissas investigated here.}
%
\begin{figure}[htbp]
    \centering
    \includegraphics[scale=0.6]{integrand_legendre_theta.pdf}
    \caption{Integrand of \cref{eq:sn_legacy_cmangl_legendre_theta_chi} for \(\rho \gg \Phi, \chi\) and \(\rho \ll \Phi, \chi\).
    For \(\rho \gg \Phi, \chi\) the function is independent of \(R_0\).
    For \(\rho \ll \Phi, \chi\) the integrand divided by \(\sqrt{\rho}\) is plotted for some \(R_0\) values.}
    \label{fig:integrand-legendre-theta}
\end{figure}

\subsubsection{Convergence}
%%%%%%%%%%%%%%%%%%%%%%%%%%%
%
In \cref{fig:error-legendre-theta-1e-6}, the relative errors in \(\theta\) as calculated by \cref{eq:sn_legacy_cmangl_legendre_theta_chi_approx_plain}, \cref{eq:sn_legacy_cmangl_legendre_theta_chi_approx}, and \cref{eq:sn_legacy_cmangl_legendre_theta_g_approx} are shown as a function of impact parameter for \(\varepsilon = 10^{-4}\), \(\varepsilon = 10^{-1}\), and \(\varepsilon = 100\).
The curves obtained from \cref{eq:sn_legacy_cmangl_legendre_theta_chi_approx} (solid lines) and \cref{eq:sn_legacy_cmangl_legendre_theta_g_approx} (dotted lines) are indistinguishable.
The lines corresponding to \cref{eq:sn_legacy_cmangl_legendre_theta_chi_approx_plain} (dashed lines) show significantly larger errors for small numbers of abscissas and large impact parameters (small scattering angles), demonstrating the significance of the scaling with the numerical value of \(\pi\).
%
\begin{figure}[htbp]
    \centering
    \includegraphics[scale=0.6]{error_legendre_theta_1e-6.pdf}
    \caption{Relative errors in \(\theta\) of n-point Gauss-Legendre quadrature as a function of reduced impact parameter for three energies.
    Dashed lines: \cref{eq:sn_legacy_cmangl_legendre_theta_chi_approx_plain},
    solid lines: \cref{eq:sn_legacy_cmangl_legendre_theta_chi_approx},
    dotted lines: \cref{eq:sn_legacy_cmangl_legendre_theta_g_approx}.
    The reference simulation evaluated \cref{eq:sn_legacy_cmangl_legendre_theta_chi_approx} with 100-point Gauss-Legendre quadrature.
    The relative accuracy of the apsis is better than \(10^{-6}\).}
    \label{fig:sn_legacy_error-legendre-theta-1e-6}
\end{figure}

In \cref{fig:error-legendre-theta-1e-6}, a relative accuracy of \(10^{-6}\) was imposed on the apsis \(R_0\).
In \cref{fig:error-legendre-theta-1e-3}, the same results for the relative error are shown for a relative accuracy of \(10^{-3}\).
Now the results for \cref{eq:sn_legacy_cmangl_legendre_theta_g_approx} clearly deteriorate.
So, the saving of the evaluation of \(\Phi(R_0)\) in calculating \cref{eq:sn_legacy_cmangl_legendre_theta_g_approx} is approximately compensated for by another iteration required for obtaining \(R_0\).
Considering the risk of producing negative arguments of the square root in \cref{eq:sn_legacy_cmangl_legendre_theta_g_approx}, \cref{eq:sn_legacy_cmangl_legendre_theta_chi_approx} seems preferable.
%
\begin{figure}[htbp]
    \centering
    \includegraphics[scale=0.6]{error_legendre_theta_1e-3.pdf}
    \caption{Relative errors in \(\theta\) of n-point Gauss-Legendre quadrature as a function of reduced impact parameter for three energies.
    Dashed lines: \cref{eq:sn_legacy_cmangl_legendre_theta_chi_approx_plain},
    solid lines: \cref{eq:sn_legacy_cmangl_legendre_theta_chi_approx},
    dotted lines: \cref{eq:sn_legacy_cmangl_legendre_theta_g_approx}.
    The reference simulation evaluated \cref{eq:sn_legacy_cmangl_legendre_theta_chi_approx} with 100-point Gauss-Legendre quadrature.
    The relative accuracy of the apsis is better than  \(10^{-3}\).}
    \label{fig:sn_legacy_error-legendre-theta-1e-3}
\end{figure}
